\chapter{Estado del arte}
\label{cap:estado-arte}

Este capítulo revisa las tecnologías sobre las que se apoya ComprendiA y las alternativas
consideradas en cada caso, justificando las decisiones tomadas.

\section{Transcripción automática del habla}

La conversión de voz a texto (\textit{Automatic Speech Recognition}, ASR) es el primer paso del
sistema. Se valoraron dos vías:

\begin{itemize}
  \item \textbf{Subtítulos de YouTube.} Muchos vídeos disponen de subtítulos, ya sean manuales o
        generados automáticamente por la plataforma. Reutilizarlos evita descargar y procesar el
        audio, reduciendo notablemente el tiempo y el coste.
  \item \textbf{Whisper} \cite{radford2022whisper}. Modelo de reconocimiento del habla de OpenAI,
        robusto frente a ruido y con buen rendimiento multilingüe. Se emplea como alternativa
        cuando el vídeo no ofrece subtítulos válidos.
\end{itemize}

ComprendiA adopta una estrategia híbrida: intenta primero obtener los subtítulos mediante
\texttt{yt-dlp} y solo recurre a Whisper si no existen o son de calidad insuficiente. Esta
decisión prioriza el rendimiento sin renunciar a la cobertura.

\section{Representación semántica: \textit{embeddings}}

Para poder buscar por significado y no solo por coincidencia de palabras, el texto se
transforma en vectores numéricos (\textit{embeddings}) que sitúan fragmentos de contenido
similar cerca en el espacio vectorial. Se utiliza el modelo \texttt{text-embedding-3-small} de
OpenAI por su equilibrio entre calidad y coste. La similitud entre vectores se mide con la
distancia del coseno.

El almacenamiento y la consulta eficiente de estos vectores se resuelven con la extensión
\texttt{pgvector} \cite{pgvector} de PostgreSQL, que añade un tipo de dato vectorial y
operadores de distancia indexables a una base de datos relacional convencional, evitando
introducir una base de datos vectorial dedicada.

\section{Búsqueda semántica y generación aumentada por recuperación (RAG)}

La arquitectura RAG \cite{lewis2020rag} combina la recuperación de información con la generación
de lenguaje natural: ante una pregunta, primero se recuperan los fragmentos más relevantes
mediante búsqueda semántica y después se entregan a un modelo de lenguaje como contexto para que
redacte la respuesta. Frente al uso directo de un modelo generativo, RAG reduce las alucinaciones
y ancla las respuestas en el contenido real del vídeo, además de permitir citar el momento exacto
de la grabación.

\section{Modelos de lenguaje generativos}

Para la redacción de respuestas, la generación de resúmenes y la extracción de capítulos y
conceptos se emplea \texttt{gpt-4o-mini} de OpenAI, un modelo de coste contenido y latencia
razonable, suficiente para las tareas planteadas. El sistema está diseñado para degradar de
forma controlada: si el modelo no está disponible, se aplican alternativas locales (por ejemplo,
una segmentación por bloques temporales) para no interrumpir el procesamiento.

\section{Herramientas y plataformas similares}

% TODO: ampliar con un breve estudio comparativo (p. ej. funciones de resumen de vídeo de
% terceros, extensiones de navegador, plataformas LMS) y una tabla comparativa.
Existen soluciones que abordan partes del problema —generadores de resúmenes de vídeo,
extensiones que transcriben YouTube o asistentes de estudio genéricos—, pero no integran en un
único flujo la transcripción, el análisis estructurado (capítulos y conceptos), la consulta
conversacional anclada en el vídeo y la organización por asignaturas. ComprendiA se diferencia
precisamente por esa integración orientada al estudio.
